\section{Materials}

\subsection{Equipment}
Hardware:

\begin{itemize}
  \item Minimum configuration: 16 GB RAM, 8 CPU cores, 50 GB available storage
  \item Recommended configuration: 32 GB RAM, 16 CPU cores, 100 GB available storage
  \item High-performance setup: 64 GB RAM, 32 CPU cores, 250 GB SSD storage
  \item Operating systems: Linux (Ubuntu 20.04+), macOS (12.0+), and Windows 10/11 (with Windows Subsystem for Linux recommended)
\end{itemize}

\subsection{Software dependencies}
Core requirements are:

\begin{itemize}
  \item Python 3.12 or higher with pip package manager
  \item Conda package manager (Miniconda or Anaconda)
  \item R statistical computing environment (version 4.1+)
\end{itemize}
Python packages:

\begin{itemize}
  \item cutadapt (4.9+): Sequence adapter trimming and demultiplexing
  \item rdkit (2024.3+): Chemical structure processing and property calculation
  \item pandas (2.2+): Data manipulation and statistical analysis
  \item numpy (1.24+): Numerical computing and array operations
  \item matplotlib/seaborn: Data visualization and publication-quality plotting
  \item plotly: Interactive visualization and dashboard components
  \item scipy (1.10+): Statistical functions and optimization algorithms
\end{itemize}
R packages:

\begin{itemize}
  \item edgeR: Differential expression analysis adapted for count data
  \item limma: Linear modeling and empirical Bayes statistics
  \item tidyverse: Data manipulation and visualization tools
  \item GGally: Extension to ggplot2 for correlation matrices
  \item BiocManager: Bioconductor package management
\end{itemize}

\subsection{Data}
We provide a complete single-display DELT-Hit example dataset to demonstrate the full workflow from initialization through demultiplexing, library reconstruction, and enrichment analysis. A complementary dual-display example focused on library enumeration and visualization is also included in the supporting material. The example is based on the architecture of a published library\cite{bassi_singlestrandeddnaencodedchemical_2020} and can be accessed through the repository's supporting material.
The dataset includes \texttt{example-single-display.xlsx}, which defines the library architecture and experimental setup; \texttt{campaign.fastq.gz}, which contains representative sequencing reads from the screen; and \texttt{analysis.yaml}, which specifies the comparisons used for hit identification. The files \texttt{example-single-display.xlsx} and \texttt{analysis.yaml} are available in the GitHub repository under \path{supporting_material/experiments/example-single-display/}. Table~\ref{tab:table2} lists the tutorial inputs and persistent archives. A complete end-to-end walkthrough of this example is provided in Box~2 in the Procedure.

\begin{table}[htbp]
\centering
\caption{Tutorial datasets for protocol validation}
\label{tab:table2}
\begin{tabular}{p{4.0cm}p{6.5cm}p{4.0cm}}
\toprule
File & URL & Description \\
\midrule
campaign.fastq.gz & \url{https://ndownloader.figshare.com/files/61487743} & Representative screening data for the single-stranded example \\
example-single-stranded.xlsx & \path{supporting_material/experiments/example-single-stranded/} & Complete library and experimental definitions for the full single-stranded workflow \\
analysis.yaml & \path{supporting_material/experiments/example-single-stranded/} & Experiment configuration for enrichment analysis \\
Full workflow archive & \url{https://doi.org/10.5281/zenodo.20447074} & Persistent archive containing the single-stranded workflow data together with generated outputs \\
Code release archive & \url{https://doi.org/10.5281/zenodo.20497548} & Persistent archive of the DELT-Hit code release used to produce the workflow outputs \\
\bottomrule
\end{tabular}
\end{table}


\clearpage
\begin{tcolorbox}[
  colback=gray!10,
  colframe=gray!50,
  title=Box 1 | Installing DELT-Hit,
  fonttitle=\bfseries
]

\setupitem{Environment setup with Conda}

We recommend using Miniconda to create an isolated environment with controlled dependencies. This step is only required once.


\begin{lstlisting}
# Download and install Miniconda for your operating system
# Follow instructions at: https://docs.anaconda.com/miniconda

# Create dedicated environment (only needed once)
conda create -n delt-hit python=3.12 -y

# Activate environment (required for each new terminal session)
conda activate delt-hit
conda install pygraphviz -y

# Install DELT-Hit package
pip install git+https://git@github.com/DELTechnology/delt-hit.git
\end{lstlisting}
<CRITICAL>: Activate the \texttt{delt-hit} conda environment using \texttt{conda activate delt-hit} each time you open a new terminal session before running DELT-Hit commands. We also recommend initializing the shell with Conda during installation so that the \texttt{conda} command is available in new terminal windows. 


\setupitem{R environment configuration}

Statistical analysis requires R and specific Bioconductor packages. This configuration is only required once. Open R or RStudio and execute:


\begin{lstlisting}
# Install required CRAN packages
install.packages(c("tidyverse", "GGally"))

# Install BiocManager for Bioconductor packages
if (!require("BiocManager", quietly = TRUE))
    install.packages("BiocManager")

# Install Bioconductor packages
BiocManager::install(c("edgeR", "limma"))
\end{lstlisting}

\setupitem{Installation verification}

Verify that all components are correctly installed:


\begin{lstlisting}
# Activate environment
conda activate delt-hit

# Test core functionality
delt-hit --help

# Expected output: Help message describing the delt-hit command
\end{lstlisting}

\textbf{$\blacklozenge$ Troubleshooting}

\end{tcolorbox}
\clearpage

\subsection{Experimental setup}
The installation procedure detailed in Box~1 is performed once before running the DELT-Hit workflow. Subsequent analyses are executed from the command line within the same software environment.
